\begin{statementbox}
	\textbf{\textcolor{CStatement}{条件}}
	\begin{description}[style=nextline, leftmargin=2.8em, labelsep=0.5em, itemsep=0.35em]
		\item[\textbf{\textcolor{CStatement}{条件1（定义限制）：}}]
			$x \neq \pi + 2k\pi$ ($k \in \mathbb{Z}$)，确保 $\tan\frac{x}{2}$ 有定义。
		\item[\textbf{\textcolor{CStatement}{条件2（分母非零）：}}]
			对 $\tan x$ 公式还需 $x \neq \frac{\pi}{2} + k\pi$ ($k \in \mathbb{Z}$)，保证分母 $1-t^2 \neq 0$。
	\end{description}

	\medskip
	\textbf{\textcolor{CStatement}{结论}}
	\begin{description}[style=nextline, leftmargin=2.8em, labelsep=0.5em, itemsep=0.45em]
		\item[\textbf{\textcolor{CStatement}{结论一（正余弦公式）：}}]
			\textbf{\textcolor{CStatement}{直观描述：}}
			用半角正切 $t$ 统一表示 $\sin x$ 和 $\cos x$。
			\[
				\sin x = \frac{2t}{1+t^2},
				\qquad
				\cos x = \frac{1-t^2}{1+t^2}.
			\]
		\item[\textbf{\textcolor{CStatement}{推广结论（正切公式）：}}]
			\textbf{\textcolor{CStatement}{直观描述：}}
			在附加条件下，$\tan x$ 也可用 $t$ 表示。
			\[
				\tan x = \frac{2t}{1-t^2}.
			\]
	\end{description}

\end{statementbox}
